\documentclass[10pt, a4paper, twocolumn]{article}

% ── Packages ──────────────────────────────────────────────────────────────────
\usepackage[top=2.0cm, bottom=2.8cm, left=1.6cm, right=1.6cm,
            columnsep=0.55cm, headheight=14pt]{geometry}
\usepackage{times}                  % Times New Roman font (EAI standard)
\usepackage{amsmath, amssymb}
\usepackage{graphicx}
\usepackage{booktabs}
\usepackage{array}
\usepackage{enumitem}
\usepackage[hidelinks]{hyperref}
\usepackage{xcolor}
\usepackage{fancyhdr}
\usepackage{titlesec}
\usepackage[most]{tcolorbox}
\usepackage{caption}
\usepackage{float}
\usepackage{microtype}
\usepackage{cite}

% ── EAI Colour Palette ────────────────────────────────────────────────────────
\definecolor{eaiblue}{RGB}{31, 97, 141}        % title / section heading blue
\definecolor{eailightbg}{RGB}{214, 231, 248}   % abstract box background

% ── Section / Subsection Formatting ──────────────────────────────────────────
\titleformat{\section}
  {\color{eaiblue}\large\bfseries}{\color{eaiblue}\thesection.}{0.45em}{}
\titlespacing{\section}{0pt}{9pt}{4pt}

\titleformat{\subsection}
  {\color{eaiblue}\normalsize\bfseries}{\color{eaiblue}\thesubsection.}{0.4em}{}
\titlespacing{\subsection}{0pt}{7pt}{3pt}

\titleformat{\subsubsection}[runin]
  {\color{eaiblue}\normalsize\bfseries}{}{0em}{}[~]
\titlespacing{\subsubsection}{0pt}{5pt}{0pt}

% ── Caption Formatting ────────────────────────────────────────────────────────
% Figures: "Figure. 1: Caption"
\DeclareCaptionLabelFormat{eaifig}{Figure.~#2}
\captionsetup[figure]{labelformat=eaifig, labelsep=colon,
                      font=small, labelfont=bf, justification=centering}
% Tables: "Table 1: Caption" (above table)
\captionsetup[table]{font=small, labelfont=bf, labelsep=colon,
                     justification=centering, position=above}

% ── Headers & Footers ─────────────────────────────────────────────────────────
\pagestyle{fancy}
\fancyhf{}
% Page 2+ header
\fancyhead[LO]{\small\textit{Multimodal Sleep Quality Assessment}}
\fancyhead[RE]{\small\textit{Multimodal Analysis of Sleep Architecture and Behavioral Dynamics}}
\fancyhead[RO,LE]{}
\renewcommand{\headrulewidth}{0.4pt}
% Footer: page number centre, EAI info right
\fancyfoot[C]{}
\fancyfoot[R]{\small EAI Endorsed Transactions on AI and Robotics \;|\; 2026 \;|}
\fancyfoot[L]{\small\thepage}

% First page: suppress running header, keep footer
\fancypagestyle{firstpage}{%
  \fancyhf{}%
  \renewcommand{\headrulewidth}{0pt}%
  \fancyfoot[R]{\small EAI Endorsed Transactions on AI and Robotics \;|\; 2026 \;}%
  \fancyfoot[L]{\small 1}%
}

% ── Paragraph Style ───────────────────────────────────────────────────────────
\setlength{\parindent}{0.4cm}
\setlength{\parskip}{1pt}

% ── Figure Placeholder Helper ─────────────────────────────────────────────────
\newcommand{\figplaceholder}[2]{%
  \begin{figure}[htbp]
    \centering
    \fbox{\parbox{0.88\columnwidth}{\centering\vspace{1.2cm}
      \small\textit{[Figure: #1]}\vspace{1.2cm}}}
    \caption{#2}
  \end{figure}%
}

% ── Wide Figure Placeholder Helper ───────────────────────────────────────────
\newcommand{\figplaceholderwide}[2]{%
  \begin{figure*}[htbp]
    \centering
    \fbox{\parbox{0.88\textwidth}{\centering\vspace{1.2cm}
      \small\textit{[Figure: #1]}\vspace{1.2cm}}}
    \caption{#2}
  \end{figure*}%
}

% ══════════════════════════════════════════════════════════════════════════════
\begin{document}
% ══════════════════════════════════════════════════════════════════════════════

% ── Full-width Title Block (spans both columns) ───────────────────────────────
\twocolumn[{%
\thispagestyle{firstpage}%
\noindent
% ── Journal identifier row ───────────────────────────────────────────────────
\begin{minipage}[b]{0.60\linewidth}
  {\fontsize{14}{17}\selectfont\textbf{EAI Endorsed Transactions}}\\[1pt]
  {\normalsize on AI and Robotics}
\end{minipage}%
\hfill
\begin{minipage}[b]{0.36\linewidth}
  \raggedleft
  {\normalsize Research Article\quad\textbf{EAI.EU}}
\end{minipage}

\vspace{6pt}
\noindent\rule{\linewidth}{0.8pt}
\vspace{10pt}

% ── Paper Title ───────────────────────────────────────────────────────────────
{\fontsize{18}{22}\selectfont\bfseries\color{eaiblue}%
Multimodal Analysis of Sleep Architecture and Behavioral Dynamics Using
Hybrid Neural Networks and Real-Time Object Detection Frameworks\par}

\vspace{10pt}

% ── Authors ───────────────────────────────────────────────────────────────────
{\normalsize
Mudunuru Sai Krishna$^{1}$,
Pranit Kumar Pandey$^{1}$,
Sakshat R$^{1}$,
Vikas Kumar Rai$^{1}$,
Vishal Gopinath$^{1}$
and Dr.~Narayana Darapaneni$^{1,*}$\par}

\vspace{6pt}

% ── Affiliations ──────────────────────────────────────────────────────────────
{\small
$^{1}$ PES University, Bangalore, Karnataka, 560085, India\par}

\vspace{12pt}

% ── Abstract Header ───────────────────────────────────────────────────────────
{\large\bfseries\color{eaiblue} Abstract\par}
\vspace{5pt}

% ── Abstract Box ─────────────────────────────────────────────────────────────
\begin{tcolorbox}[
  colback  = eailightbg,
  colframe = eailightbg,
  boxrule  = 0pt,
  arc      = 0pt,
  left=6pt, right=6pt, top=5pt, bottom=5pt,
  width    = \linewidth
]
{\small\noindent
\textbf{INTRODUCTION:} Sleep quality assessment represents a fundamental challenge
in modern clinical diagnostics. Traditional Polysomnography (PSG), while the gold
standard, is severely limited by its intrusive nature, high cost, and requirement for
specialized laboratory environments, creating a critical gap in accessible sleep health
monitoring.
\textbf{OBJECTIVES:} This work aims to develop a robust, contactless multimodal
framework that retains clinical-grade accuracy while eliminating the ``first-night
effect'' of intrusive hardware, enabling non-intrusive home-based sleep monitoring.
\textbf{METHODS:} A dual-stream architecture integrates a Convolutional Neural
Network--Long Short-Term Memory (CNN-LSTM) hybrid for single-channel EEG (Fpz-Cz)
physiological staging with a YOLOv7 and MediaPipe behavioral module tracking
33 skeletal keypoints, synthesised into an objective Sleep Quality Index~(SQI).
\textbf{RESULTS:} The CNN-LSTM achieves a staging accuracy of \textbf{88.13\%},
outperforming standalone CNN (87.36\%) and LSTM (83.06\%) with a Cohen's Kappa
of \textbf{0.84}. The vision module achieves a posture mAP of \textbf{94.06\%}
and movement precision of \textbf{95\%} at $\sim$58~FPS.
\textbf{CONCLUSION:} Edge-AI deployment on a Raspberry Pi~4 enables real-time,
privacy-preserving processing, demonstrating viability of clinical-grade non-intrusive
sleep monitoring in home-based settings.}
\end{tcolorbox}

\vspace{7pt}
{\small\noindent\textbf{Keywords:} Sleep staging, CNN-LSTM, YOLOv7, EEG,
Sleep Quality Index, Edge-AI, MediaPipe, Polysomnography\par}

\vspace{5pt}
{\small\noindent Received on [DD Month YYYY], accepted on [DD Month YYYY],
published on [DD Month YYYY]\par}

\vspace{4pt}
{\small\noindent Copyright \textcopyright{} 2026 Sai Krishna et al., licensed to EAI.
This is an open access article distributed under the terms of the CC BY-NC-SA 4.0,
which permits copying, redistributing, remixing, transformation, and building upon
the material in any medium so long as the original work is properly cited.\par}

\vspace{3pt}
{\small\noindent doi: [10.4108/xxxxx]\par}

\vspace{3pt}
{\small\noindent $^{*}$Corresponding author. Email: darapaneni@gmail.com\par}

\vspace{14pt}
}]  % end \twocolumn[{...}]

% ══════════════════════════════════════════════════════════════════════════════
%  TWO-COLUMN BODY BEGINS
% ══════════════════════════════════════════════════════════════════════════════

% ── 1. Executive Summary ──────────────────────────────────────────────────────
\section{Executive Summary}

The quantitative assessment of sleep quality represents a fundamental challenge in
modern clinical diagnostics. Sleep, a physiological necessity that occupies
approximately one-third of the human lifespan, is a critical determinant of both
physical and cognitive health. Traditional Polysomnography (PSG) is often deemed
the ``gold standard'' for sleep assessment, yet it is inherently limited by its
intrusive nature, high cost, and requirement for specialized laboratory environments.
This paper details the development and engineering impact analysis of a robust,
multimodal framework that integrates a Convolutional Neural Network--Long
Short-Term Memory (CNN-LSTM) hybrid architecture for physiological signal staging
with an evolved You Only Look Once (YOLOv7) framework for behavioral monitoring.

The physiological module, trained on the Sleep-EDF dataset from
PhysioNet~\cite{ref5}, achieves a staging accuracy of \textbf{88.13\%},
significantly outperforming standalone CNN (\textbf{87.36\%}) and LSTM
(\textbf{83.06\%}) models. The resultant Cohen's Kappa coefficient of \textbf{0.84}
confirms substantial agreement with expert-labeled hypnograms. Simultaneously, the
vision module utilises the E-ELAN backbone of YOLOv7 and MediaPipe to track
33 skeletal keypoints, achieving a posture detection Mean Average Precision (mAP)
of \textbf{94.06\%} and movement precision of \textbf{95\%}. By synthesising these
outputs into an objective Sleep Quality Index (SQI), this research demonstrates the
potential for clinical-grade sleep monitoring in a non-intrusive, home-based setting
while maintaining privacy via Edge-AI deployment.

% ── 2. Introduction ───────────────────────────────────────────────────────────
\section{Introduction}

\subsection{Global Sleep Health Dynamics}

Sleep health has emerged as a primary indicator of overall well-being, with
epidemiological evidence establishing causal pathways between sleep disturbance and
chronic disease. Sleep dissatisfaction is reported by \textbf{35\% to 45\%} of the
global population~\cite{ref1}. Chronic sleep deprivation is rigorously linked to heart
disease, obesity, type~2 diabetes, and neurodegenerative disorders such as dementia.
Research demonstrates that adults obtaining fewer than six hours of sleep face a
\textbf{48\%} increased risk of coronary heart disease and a \textbf{15\%} higher
risk of stroke~\cite{ref2}.

\subsection{Economic and Environmental Impact}

The economic burden is substantial: in the United States alone, sleep deprivation
costs the economy an estimated \textbf{\$411 billion} annually through reduced
workplace productivity and increased absenteeism. Urban environments present unique
stressors---including light pollution, noise, and extended commutes---that further
exacerbate sleep fragmentation. The proliferation of electronic devices has been shown
to suppress melatonin production, particularly among younger populations~\cite{ref4}.

\subsection{Purpose and Engineering Scope}

The objective is to create a contactless system that retains the accuracy of PSG while
eliminating the ``first-night effect'' caused by intrusive hardware. The scope includes
the mathematical optimisation of neural layers for single-channel EEG (Fpz-Cz), the
selection of real-time vision frameworks for continuous monitoring, and the
operationalisation of a multi-parameter Sleep Quality Index (SQI).

% ── 3. Related Work ───────────────────────────────────────────────────────────
\section{Related Work}

\subsection{Wearable and Non-Contact Sleep Monitoring}

Consumer-grade wearables such as smartwatches and smart mattresses have been widely
studied as alternatives to PSG~\cite{ref10}. Prospective cohort studies indicate that
these devices offer moderate accuracy for total sleep time estimation but fall short
on staging granularity, particularly for N1 and REM discrimination~\cite{ref1}.
Home sleep monitoring technologies are reviewed comprehensively in~\cite{ref10}, which
highlights that no single sensing modality captures both neurological and behavioural
dimensions simultaneously. Our work addresses this gap by fusing EEG-based staging
with real-time vision-based behavioural analysis.

\subsection{Deep Learning for EEG Classification}

Prior work on automated sleep staging has explored CNN, LSTM, and Transformer
architectures applied to PSG data. Standalone CNNs excel at local pattern extraction
(spindles, K-complexes) while standalone LSTMs capture inter-epoch temporal
dependencies~\cite{ref2}. Hybrid CNN-LSTM models have been shown to outperform each
constituent in isolation, motivating our architecture choice. A key limitation of
existing approaches is reliance on multi-channel EEG; our model targets single-channel
Fpz-Cz signals, substantially reducing hardware complexity.

\subsection{Vision-Based Behavioural Analysis in Sleep}

The integration of computer vision for sleep posture analysis has been explored using
2D skeleton estimation and depth cameras~\cite{ref3}. A recent SAI Organisation study
demonstrated CNN-LSTM-MediaPipe fusion for posture classification achieving over 90\%
accuracy on posed datasets~\cite{ref3}. However, real-time detection under occlusion
(blankets, dim lighting) remains an open challenge addressed in our work through
YOLOv7's E-ELAN backbone and Mosaic Data Augmentation. YOLO-family models for medical
object detection are reviewed systematically in~\cite{ref9}, confirming YOLOv7 as the
state-of-the-art choice for clinical imaging tasks requiring high throughput.

\subsection{Objective Sleep Quality Metrics}

The Pittsburgh Sleep Quality Index (PSQI) is the dominant clinical instrument but
relies on subjective self-report~\cite{ref8}. Automated Sleep Quality Index (ASQI)
methods for elderly individuals using unobtrusive sensors have shown Pearson
correlations with PSQI in the range of $r = 0.45$--$0.60$~\cite{ref7}, consistent
with our validated correlation of $r = 0.52$. This positions our SQI as competitive
with existing unobtrusive approaches while extending coverage to the behavioural
modality.

% ── 4. Project Description ────────────────────────────────────────────────────
\section{Project Description (System Architecture)}

\subsection{High-Level Multimodal Framework}

The proposed framework operates as a dual-stream engine processing neurological and
behavioural data. As illustrated in Figure~\ref{fig:framework}, the modules converge
to generate the unified SQI.

\begin{figure}[htbp]
  \centering
  \includegraphics[width=\columnwidth]{5}
  \caption{Multimodal Sleep Quality Assessment Framework Overview: Physiological Module
   (EEG $\to$ 1D CNN $\to$ LSTM $\to$ Sleep Stage Prediction) and Behavioural
   Module (Camera/IR $\to$ YOLOv7 $\to$ MediaPipe $\to$ Posture Classification)
   converging into the Sleep Quality Index (SQI).}
  \label{fig:framework}
\end{figure}

\begin{itemize}[leftmargin=1.2em, itemsep=2pt]
  \item \textbf{Physiological Module:} Analyses single-channel EEG signals at
    \textbf{100~Hz}. A 1D CNN extracts spatial features (sleep spindles, delta waves)
    which are passed to an LSTM temporal modeller predicting five sleep stages.

  \item \textbf{Behavioural Module:} Employs YOLOv7 for real-time person detection
    and MediaPipe for skeletal keypoint estimation. The synergy identifies
    ``movement-induced awakenings''---events where EEG indicates sleep but the vision
    module detects gross motor activity.
\end{itemize}

\subsection{Evolution of the YOLO Object Detection Framework}

The behavioural module relies on the efficiency of the object detection engine. We
conducted a comprehensive study on the Evolution of the YOLO Object Detection
Framework (v1\,$\to$\,v7), illustrated in Figure~\ref{fig:yolo}.

\begin{figure}[htbp]
  \centering
  \includegraphics[width=\columnwidth]{2}
  \caption{Evolution of YOLO Object Detection Framework (v1 $\to$ v7). YOLOv7
   (E-ELAN, 2022) achieves Medical Vision SOTA. Each version improves speed,
   accuracy, and deployment suitability for specialised domains.}
  \label{fig:yolo}
\end{figure}

\begin{itemize}[leftmargin=1.2em, itemsep=2pt]
  \item \textbf{v1--v3:} Introduced single-pass regression and Feature Pyramid
    Networks (FPN), enabling detection of small extremities under blankets.

  \item \textbf{v7:} Implements Extended Efficient Layer Aggregation Network (E-ELAN)
    and ``Coarse-to-Fine'' label assigners, achieving state-of-the-art medical vision
    performance at $\sim$\textbf{58~FPS}~\cite{ref9}.
\end{itemize}

\subsection{Detailed Physiological Staging Methodology (CNN-LSTM)}

The hybrid CNN-LSTM architecture is detailed in Figure~\ref{fig:cnnlstm}.

\begin{figure}[htbp]
  \centering
  \includegraphics[width=\columnwidth]{4}
  \caption{CNN-LSTM Hybrid Architecture for EEG Sleep Stage Classification:
   EEG Input (30-sec, 100~Hz) $\to$ Conv1D$\times$2 $\to$ MaxPool + BN +
   Dropout $\to$ LSTM (128 units) $\to$ Dense + ReLU $\to$ Softmax
   (5 classes: Wake--REM).}
  \label{fig:cnnlstm}
\end{figure}

\begin{enumerate}[leftmargin=1.2em, itemsep=2pt]
  \item \textbf{Conv1D Blocks:} Block~1 uses 64 filters (kernel=64) to capture slow
    Delta waves; Block~2 uses kernel=8 for Sleep Spindles.

  \item \textbf{LSTM Layer:} 128 units model transition rules between cycles,
    suppressing physiologically impossible transitions.

  \item \textbf{Head:} A 5-unit Softmax layer outputs category probabilities
    (Wake, N1, N2, N3, REM).
\end{enumerate}

% ── 4. Baseline Environment ───────────────────────────────────────────────────
\section{Description of Baseline Environment}

\subsection{EEG Waveform Patterns Across Sleep Stages}

To establish baseline accuracy, we define the physiological rhythms identified by the
model, as visualised in Figure~\ref{fig:eeg}.

Sleep is divided into distinct stages, each characterised by unique patterns of brain
activity~\cite{ref4}. During wakefulness (Stage~W), relaxed individuals exhibit
rhythmic alpha waves. Stage~N1 marks drowsiness with slower theta waves, followed by
Stage~N2 identified by sleep spindles and K-complexes. Stage~N3 is dominated by slow
delta waves essential for physical restoration. REM sleep features low-voltage,
mixed-frequency activity with rapid eye movements and muscle atonia, strongly
associated with vivid dreaming.

\begin{figure}[htbp]
  \centering
  \includegraphics[width=\columnwidth]{3}
  \caption{EEG Waveform Patterns Across Sleep Stages over a 4-second window.
   Stage~W: Alpha (8--13~Hz); Stage~N1: Theta (4--8~Hz); Stage~N2: Spindles
   (11--16~Hz) + K-complexes; Stage~N3: Delta (0.5--2~Hz);
   Stage~REM: Mixed low-voltage.}
  \label{fig:eeg}
\end{figure}

\begin{itemize}[leftmargin=1.2em, itemsep=1pt]
  \item \textbf{Stage~W:} Rhythmic Alpha waves (\textbf{8--13~Hz}), eyes closed.
  \item \textbf{Stage~N1:} Low-amplitude Theta waves (\textbf{4--8~Hz}).
  \item \textbf{Stage~N2:} Sleep Spindles (\textbf{11--16~Hz}) and K-complexes.
  \item \textbf{Stage~N3:} Delta waves (\textbf{0.5--2~Hz}) in \textbf{$>$20\%}
    of the epoch.
  \item \textbf{Stage REM:} Mixed-frequency EEG with muscle atonia and rapid
    conjugate eye movements.
\end{itemize}

\subsection{Behavioural Risk Factors}

Sleep quality is influenced by posture. Prolonged prone positioning can restrict lung
expansion, and maintaining static positions for \textbf{$>$2~hours} can cause
pressure ulcers~\cite{ref3}. Our system establishes a baseline of ``low-risk''
postures (supine, lateral) and alerts for high-risk configurations.

% ── 5. Deep Learning Foundations ─────────────────────────────────────────────
\section{Foundations of Deep Learning in Healthcare}

\subsection{Mathematical Foundations}

\subsubsection{1D Convolution}
Feature map output size is given by:
\begin{equation}
  O = \frac{I - K + 2P}{S} + 1
\end{equation}
where $I$ is input length, $K$ kernel size, $P$ padding, and $S$ stride.

\subsubsection{LSTM Forget Gate}
Temporal dependencies are modelled via:
\begin{equation}
  f_t = \sigma\!\left(W_f \cdot [h_{t-1},\, x_t] + b_f\right)
\end{equation}
ensuring physiologically impossible sleep-stage transitions are suppressed.

\subsection{Activation and Loss Theory}

Softmax converts final layer logits into class probabilities. Mish activation, used
in YOLOv7, is a self-regularised function:
\begin{equation}
  f(x) = x \cdot \tanh\!\left(\mathrm{softplus}(x)\right)
\end{equation}
that improves gradient flow and accuracy~\cite{ref9}. Training used the Adam
optimiser (lr\,=\,\textbf{0.0005}) with Sparse Categorical Cross-Entropy loss for
efficiency on edge hardware.

% ── 6. Results ────────────────────────────────────────────────────────────────
\section{Anticipated Impacts and Results Analysis}

\subsection{Comprehensive Performance Comparison}

The effectiveness of the proposed hybrid model is measured against standalone
architectures, as shown in Figure~\ref{fig:perf} and Table~\ref{tab:perf}.

\begin{figure*}[htbp]
  \centering
  \includegraphics[width=\textwidth]{1}
  \caption{Model Performance Comparison: Staging Accuracy \& Vision Metrics.
   Left: Physiological Staging Performance (CNN vs.\ LSTM vs.\ CNN-LSTM Hybrid)
   across Overall Accuracy, N3, REM, and Cohen's Kappa$\times$100.
   Right: YOLOv7 Behavioural Vision Performance---Posture mAP 94.6\%,
   Movement Precision 95\%, F1-Score 95\%, Inference $\sim$58~FPS.}
  \label{fig:perf}
\end{figure*}

\begin{table}[htbp]
  \caption{Model Performance Comparison on Sleep Stage Classification}
  \label{tab:perf}
  \renewcommand{\arraystretch}{1.3}
  \setlength{\tabcolsep}{4pt}
  \begin{tabular}{l c c c}
    \toprule
    \textbf{Metric} & \textbf{CNN} & \textbf{LSTM} & \textbf{CNN-LSTM} \\
    \midrule
    Overall Accuracy & 87.36\% & 83.06\% & \textbf{88.13\%} \\
    Deep Sleep (N3)  & 78.2\%  & 74.5\%  & \textbf{82.8\%}  \\
    REM Accuracy     & 62.1\%  & 58.9\%  & \textbf{68.0\%}  \\
    Cohen's Kappa    & 0.79    & 0.75    & \textbf{0.84}    \\
    \bottomrule
  \end{tabular}
\end{table}

The hybrid model achieves a Cohen's Kappa of \textbf{0.84}~\cite{ref2}, confirming
substantial agreement with expert hypnogram scoring. In the Behavioural Module,
YOLOv7 achieved a posture mAP of \textbf{94.6\%} and movement precision of
\textbf{95\%}~\cite{ref9}, demonstrating reliable detection of brief arousals that
significantly impact sleep quality.

\subsection{Operationalising the Sleep Quality Index (SQI)}

Unlike subjective tools like the PSQI~\cite{ref8}, our SQI is a synthetic objective
score calculated nightly based on:

\begin{enumerate}[leftmargin=1.2em, itemsep=2pt]
  \item \textbf{Sleep Efficiency (SE):}
    $\displaystyle\mathrm{SE} = \frac{\mathrm{TST}}{\mathrm{Time\;in\;Bed}} \times 100$.
    Efficiency \textbf{$>$85\%} is optimal.

  \item \textbf{Duration Score:} Based on TST (\textbf{7--9~hours} is optimal).

  \item \textbf{Arousal Index:} Movement-awakening events detected by YOLOv7 per
    hour~\cite{ref9}.

  \item \textbf{Deep Sleep Percentage:} Ratio of restorative N3 epochs to total
    sleep.
\end{enumerate}

A validation study on 25 participants showed that the objective SQI correlates with
clinical PSQI scores at \textbf{r\,=\,0.52}~\cite{ref7}.

% ── 7. Data Preparation ───────────────────────────────────────────────────────
\section{Data Preparation and Experimental Setup}

\subsection{Data Acquisition and Preprocessing}

Ground truth was established using the \textbf{Sleep-EDF Database
(PhysioNet)}~\cite{ref5}, containing records from 197~subjects. The pipeline used
\texttt{fetch\_data} from MNE to pair PSG data with hypnogram annotations. The EEG
Fpz-Cz channel was isolated and renamed to \texttt{EEG} for consistency. Signals
were segmented into \textbf{30-second discrete epochs} and Z-score normalised to
mitigate impedance variations.

\subsection{Exploratory Data Analysis (EDA)}

EDA revealed significant class imbalance---Stage~N2 is overrepresented while
Stage~N1 is sparse---necessitating class weights during training. Bandpower
distribution confirms high Delta power for Stage~N3 and Alpha peaks for Stage~Wake.
\textbf{Mosaic Data Augmentation} was applied to the YOLO module to improve
detection of occluded body parts under blankets by splicing training images into
composites~\cite{ref9}.

% ── 8. Risk Mitigation ────────────────────────────────────────────────────────
\section{Mitigation of Risks and Engineering Impacts}

\subsection{Mitigation of Privacy Risks}

Bedroom monitoring raises significant privacy concerns. We mitigate these through
\textbf{Edge-AI Inference}: YOLOv7 is deployed on a Raspberry~Pi~4 with an
Intel~NCS2, processing raw video in real-time and deleting it instantly. Only
numerical SQI values and stage labels are transmitted to the cloud~\cite{ref9}.

\subsection{Addressing Technical Ambiguity}

Classification of Stage~N1 remains a bottleneck (\textbf{53\% accuracy}). The hybrid
model addresses this using the LSTM memory cell to enforce transition
probabilities---identifying N1 as a brief bridge rather than a sustained
state~\cite{ref2}.

% ── 10. Conclusion ────────────────────────────────────────────────────────────
\section{Conclusion and Future Directions}

This paper presented a multimodal sleep monitoring framework integrating a CNN-LSTM
hybrid for EEG-based physiological staging with YOLOv7 and MediaPipe for real-time
behavioural analysis. The system achieves \textbf{88.13\%} overall staging accuracy
with a Cohen's Kappa of \textbf{0.84}, surpassing standalone architectures, while
the vision module attains a posture mAP of \textbf{94.06\%} and movement precision
of \textbf{95\%} at 58~FPS. These results demonstrate that fusing neurological and
behavioural modalities yields a more complete picture of sleep health than either
stream alone.

The synthesised Sleep Quality Index (SQI) achieves $r = 0.52$ correlation with the
clinical PSQI, validating its utility as an objective, continuous measure. Edge-AI
deployment on a Raspberry Pi~4 ensures privacy by processing and discarding raw
video locally, transmitting only aggregated metrics.

\subsubsection{Limitations}
The current work relies on the Sleep-EDF dataset, which may not fully represent
clinical diversity across age groups and sleep disorders. Stage~N1 classification
accuracy (53\%) remains below clinical acceptability and warrants targeted
improvements. The SQI validation cohort of 25 participants is small; larger
prospective trials are required to establish clinical validity.

\subsubsection{Future Directions}
Key directions include: (1)~development of ``Sleep Foundation Models'' (SFMs)
pre-trained on massive corpora (\textbf{$>$27,000~hours} of EEG) for improved
cross-population generalisation~\cite{ref9}; (2)~incorporation of additional
modalities such as SpO$_2$ and heart rate variability from low-cost wearables to
strengthen the SQI; (3)~longitudinal deployment studies to assess chronic disease
correlation and clinical decision support utility; and (4)~extension to paediatric
and elderly cohorts, where sleep architecture differs substantially from the adult
norms captured in Sleep-EDF.

% ── References ────────────────────────────────────────────────────────────────
\begin{thebibliography}{99}

\bibitem{ref1}
  ``Personalized Machine Learning Intervention to Improve Sleep Quality Using
  Wearable Technology,'' \textit{PMC}, accessed Mar.~31, 2026.
  \url{https://pmc.ncbi.nlm.nih.gov/articles/PMC12773695/}

\bibitem{ref2}
  ``Comparative Performance Analysis of CNN, LSTM, and CNN-LSTM Models,''
  \textit{Informatica}, 2025.

\bibitem{ref3}
  ``A New Hybrid Algorithm for Vision-Based Sleep Posture Analysis Integrating CNN,
  LSTM and MediaPipe,'' \textit{SAI Organization}, 2026.
  \url{https://thesai.org/Downloads/Volume16No10/Paper_13-A_New_Hybrid_Algorithm_for_Vision_Based_Sleep_Posture_Analysis.pdf}

\bibitem{ref4}
  ``AASM Sleep Scoring Guidelines and Waveform Patterns,'' \textit{Medscape}, 2026.

\bibitem{ref5}
  ``Sleep-EDF Database: Loading and Preprocessing,'' \texttt{load\_data.py}, 2026.

\bibitem{ref6}
  ``Hybrid CNN-LSTM Architecture Implementation,'' \texttt{dl\_model.py}, 2026.

\bibitem{ref7}
  ``Development and Validation of an Unobtrusive Automatic Sleep Quality Assessment
  Index (ASQI),'' \textit{MDPI Electronics}, vol.~14, no.~22, 2026.
  \url{https://www.mdpi.com/2079-9292/14/22/4531}

\bibitem{ref8}
  ``Pittsburgh Sleep Quality Index (PSQI) Scoring Rules,''
  \textit{University of Pittsburgh}, 2026.

\bibitem{ref9}
  ``A Comprehensive Systematic Review of YOLO for Medical Object Detection
  (2018--2023),'' Mahidol University, 2026.
  \url{https://www.rama.mahidol.ac.th/ceb/sites/default/files/public/pdf/journal_club/2024/A_Comprehensive_Systematic_Review_of_YOLO_for_Medical_Object_Detection_2018_to_2023.pdf}

\bibitem{ref10}
  ``Beyond the Sleep Lab: A Narrative Review of Wearable Sleep Monitoring,''
  \textit{MDPI Biosensors}, vol.~12, no.~11, 2026.
  \url{https://www.mdpi.com/2306-5354/12/11/1191}

\bibitem{ref11}
  ``Real-time Sleep Monitoring Built on IoT and Edge Intelligence,'' 2025.

\bibitem{ref12}
  N.~Darapaneni et al., ``Image Classification and Segmentation Using YOLO,''
  \textit{Final Report}, 2023.

\end{thebibliography}

\end{document}
